\documentclass[11pt]{article}
\usepackage[margin=1in]{geometry}
\usepackage{amsmath,amssymb}
\title{E-04: Necessary Integral Stretching Blow-up}
\author{NS-BLOWUP-NECESSARY}
\date{September 3, 2026}
\begin{document}
\maketitle
\section*{1. Pregunta}
Bajo
\[
T_*<\infty,
\]
se pregunta si necesariamente
\[
\int_0^{T_*}\int_{\mathbb R^3}
\omega^T S\omega\,dx\,dt=+\infty.
\]
Se distinguen estrictamente las versiones limsup y limite.
\section*{2. Hipotesis congeladas}
Supongase una solucion de Leray--Hopf en $\mathbb R^3$ con
\[
u_0\in H^1,\qquad \operatorname{div}u_0=0.
\]
E-03 establece:
\[
\sup_{t<T_*}\|\omega(t)\|_2<\infty
\Longrightarrow T_*=\infty.
\]
Por contraposicion:
\[
T_*<\infty
\Longrightarrow
\sup_{t<T_*}\|\omega(t)\|_2=\infty.
\]
Por tanto existe una sucesion $t_n\uparrow T_*$ tal que
\[
\|\omega(t_n)\|_2\longrightarrow\infty.
\]
\section*{3. Identidad de enstrofia}
Para $0<t<T_*$, en el intervalo de regularidad fuerte,
\[
\frac12\|\omega(t)\|_2^2
+
\nu\int_0^t\|\nabla\omega(\tau)\|_2^2\,d\tau
=
\frac12\|\omega_0\|_2^2
+
\int_0^t W(\tau)\,d\tau,
\]
donde
\[
W(\tau)=
\int_{\mathbb R^3}\omega^T S\omega\,dx.
\]
Luego
\[
\int_0^t W(\tau)\,d\tau
=
\frac12\|\omega(t)\|_2^2
-\frac12\|\omega_0\|_2^2
+
\nu\int_0^t\|\nabla\omega(\tau)\|_2^2\,d\tau.
\]
En particular,
\[
\int_0^t W(\tau)\,d\tau
\ge
\frac12\|\omega(t)\|_2^2
-\frac12\|\omega_0\|_2^2.
\]
\section*{4. Resultado PROVEN: limsup}
Tomando $t=t_n$,
\[
\int_0^{t_n}W(\tau)\,d\tau\longrightarrow+\infty.
\]
Por tanto:
\[
\boxed{
\limsup_{t\uparrow T_*}
\int_0^t W(\tau)\,d\tau
=+\infty.
}
\]
Equivalentemente:
\[
\boxed{
\int_0^{T_*}W_+(\tau)\,d\tau=+\infty.
}
\]
Definimos la condicion necesaria debil:
\[
\boxed{
C_E^{\mathrm{weak}}
:=
\sup_{t<T_*}
\int_0^t
\int_{\mathbb R^3}
\omega^T S\omega\,dx\,d\tau
=+\infty.
}
\]
Por consiguiente:
\[
\boxed{
T_*<\infty
\Longrightarrow
C_E^{\mathrm{weak}}.
}
\]
\[
\boxed{E\text{-}04\ \mathrm{weak}=\mathrm{PROVEN}}
\]
\section*{5. Resultado fuerte: limite}
No se concluye de lo anterior que
\[
\lim_{t\uparrow T_*}
\int_0^tW(\tau)\,d\tau=+\infty.
\]
La funcion $W$ es firmada y la integral acumulada no es necesariamente
monotona.
Para obtener el limite seria suficiente, por ejemplo, controlar
la parte negativa:
\[
\int_0^{T_*}W_-(\tau)\,d\tau<\infty.
\]
Este control no se obtiene de E-03 por si solo.
Por tanto:
\[
\boxed{
E\text{-}04\ \mathrm{strong}=\mathrm{OPEN}.
}
\]
\section*{6. Estado ZFL}
\[
\boxed{
T_*<\infty
\Longrightarrow
C_E^{\mathrm{weak}}
\quad\mathrm{PROVEN}.
}
\]
No se identifica
\[
\int_0^{T_*}W
\]
con la version limsup.
\[
\boxed{
\text{No se usa ninguna inferencia sobre }W(t)\text{ puntual.}
}
\]
\end{document}
